\section{Conclusion}
\label{sec:conclusion}

We addressed inference-time compute allocation across the denoising trajectory of diffusion models under a strict NFE budget. Separating local noise refinement from global timestep scheduling yields a unified two-level framework that subsumes existing inference-time methods as special cases. Within this framework, \textsc{GAINS} combines offline sensitivity profiling with online feedback control. Theoretically, under a location-scale model of per-step scores, optimal offline allocation follows water-filling (\cref{thm:offline}); for adaptive policies we proved an $\Omega(T)$ regret lower bound (\cref{lem:lower}) and a matching $O(g_t^2 h\cdot T)$ upper bound for \textsc{GAINS} (\cref{thm:regret-upper}). Experiments on Stable Diffusion, EDM, and flow-based models confirm these predictions: \textsc{GAINS} reaches matched quality with 20--50\% fewer NFE, and the gains persist across local operators, prompt sets, and architectures.

\paragraph{Limitations and future work.}
The location-scale analysis assumes light-tailed per-step score distributions; heavy tails or multimodal scores would require extended analysis. Offline profiling treats sensitivities as approximately stationary across prompts; highly heterogeneous prompt distributions may benefit from prompt-conditional profiling. Beyond diffusion models, the principle of non-uniform compute allocation guided by sensitivity profiling applies wherever computation is a scarce resource in sequential decisions, including adaptive simulation budgeting and multi-stage decision problems~\citep{chen2000simulation}. The MDP formulation (\cref{sec:mdp}) admits richer policy classes (backtracking, operator mixing, multi-fidelity verifiers) that may push the quality--compute frontier further but likely require learned optimization.
